\documentclass{amsart}
\usepackage[british]{babel}
\usepackage{url}
\usepackage{hyperref}
\usepackage{xcolor}
\usepackage{enumerate}
\usepackage{tikz}
\usetikzlibrary{3d}
\usepackage{pgfplots}
\usepackage{todonotes}
\usepackage{draftwatermark}
\SetWatermarkText{\textsc{Draft -- Not for Circulation}}
\SetWatermarkScale{2}
\definecolor{edblue}{RGB}{0,0,153}
\hypersetup{
  pdftitle={\@title{}},
  colorlinks=true,
  linkcolor=edblue,anchorcolor=edblue,citecolor=edblue,filecolor=edblue,urlcolor=edblue
}%
\usepackage{geometry}
\geometry{a4paper,textwidth=5.5in,textheight=8.5in,marginparsep=7pt,marginparwidth=.6in}
\usepackage{graphicx}
\usepackage{wrapfig}
\RequirePackage[hyperref=true,
  url=true,
  isbn=true,
  backref=false,
  bibstyle=ieee,
  citestyle=authoryear-icomp,
  maxcitenames=3,
  maxbibnames=100,
  backend=biber,
  block=none]{biblatex}
\addbibresource{literature.bib}
\title{Report on new equipment for community broadband}
\author{P. Buneman \and M. Fourman \and W. Waites}
\date{\today}
\begin{document}
\begin{center}
  {\Huge\textbf{\color{red}DRAFT -- NOT FOR CIRCULATION}}
\end{center}
\maketitle
\tableofcontents
\listoffigures

\newpage

\begin{figure}[h]
  \begin{center}
    \includegraphics[width=0.49\textwidth]{sgritheall-clear-s.jpg}%
    \hfill%
    \includegraphics[width=0.49\textwidth]{sgritheall-cloudy.jpg}
  \end{center}
\label{fig-cloudy}
\caption{Scotch mist}
\end{figure}
\begin{abstract}
The Demonstrating Digital Programme of Scottish Government provided a small grant to the
University of Edinburgh to evaluate wireless (5GHz and 24GHz) and free-space optics equipment that may be useful in the
deployment of broadband in Scotland.   Ofcom
granted us an experimental licence to evaluate the
wireless equipment at
higher
power than current regulations would normally allow. The Met Office
provided access to MIDAS Land and Marine Surface Station Data. This brief report summarizes our
findings. A fuller report, with additional technical detail is
available on request.
\end{abstract}

\section{Summary}

\begin{itemize}
\item The 24 GHz equipment could prove useful over relatively short
distances (up to 8 km) if licencing permitted use at the power for
which it is designed. Under current regulations it cannot be used effectively over long
 rural links, but we are currently using it to good effect over a
 distance of a few hundred metres.

\item Given the current power limits and available bands in the 
5GHz range, the 5GHz equipment probably provides only a minor
improvement on what could be achieved by using two or three parallel
links using lower cost equipment in different bands.  Moreover the 5GHz equipment
was sensitive to tidal fading and occasionally failed completely (perhaps
due to the presence of marine radar).

The 5GHz
equipment would be much more useful if (a)  the power limits were to
be raised and (b) Ofcom were to publish their database of licences in
the C band. 

\item The free-space optics equipment performs well over
  distances of a few hundred metres in clear weather. However, it
  requires clear line-of-sight between stations. Optical
  communication fails when the visibility falls below the length of
  the link; so longer links suffer more frequent outages. 
A backup radio link is essential
  for most applications where, as across much of Scotland, conditions of low visibility may persist
  for several days. 
  \footnote{Figure~\ref{fig-cloudy} shows
    good and poor ``free-space''  visibility
    conditions on one of the shortest Tegola links.  The camera is at one relay
    and the next relay is about 2km away on the hill in the
    background.}  

The spanning-tree algorithm used to trigger fallback to radio
communications did not integrate easily with existing infrastructure,
and did not perform well in marginal conditions. We recommend that any future use of FSO equipment to supplement an
existing radio network should explore other fallback strategies that
provide better control, and include some hysteresis.

\end{itemize}


\section{Wireless Links}\label{introduction}
Many community broadband projects require a long-range link
for their connection to the Internet. Where there is no local access
to a fibre PoP, a wireless link is often the only affordable option.
The equipment of choice is
currently based on wifi operating in the unlicensed 2.4 or 5GHz
spectra. This is cheap: the equipment for a long-distance
point-to-point link costs under \pounds 400 and can provide
bi-directional throughput of about 50Mb/s. While this provides an
improvement for the many rural communities that are served by long
copper telephone lines, 50Mb/s is no longer adequate for a community
of, say, 50 residences. Technically, there is no problem in getting
more bandwidth in one of the licenced spectra, but the equipment is
more expensive and the additional cost of the licence makes this
option unaffordable for small communities.

Recently some new equipment operating at 24GHz has come on the market
from Ubiquiti\footnote{\href{http://www.ubnt.com/airfiber}{\url{http://www.ubnt.com/airfiber}}}.
This is advertised as offering 1.4Gb/s at up to 13km. Although not as
cheap (a point-to-point link costs about \pounds 3,000) it might present
an opportunity for some rural communities to upgrade their service to
be competitive with the current fibre based offerings in the UK,
assuming they can find an internet connection with that bandwidth. A
5GHz variant is also produced by Ubiquiti claiming comparable speeds
at upwards of 50km.

With this in mind, the Scottish Government's Demonstrating
Digital programme provided the University of Edinburgh with funds to
test this equipment ``in the wild''. Of course, there is ample evidence
that the equipment works, but most of the evidence we have is
from installations in urban areas over short distances. How will
it perform over longer distances in West Highland weather? And what
are the practical problems faced by communities who want to install it?

\subsection{24GHz}\label{background}

Before going into an account of our experiences, we review some
basic facts concerning wireless transmission and regulation in the 24
GHz spectrum. 

We have
already noted that the Ubiquity airFiber equipment is affordable. The advertised
throughput of 1.4Gb/s over 13+km links presumably means 700Mb/s in each direction, but
that would provide a satisfactory connection for a hundred or
so residences. Moreover, transmission in this frequency is much
less likely to be affected by tidal reflection (a significant problem
in the Highlands and Islands).

However, there are significant obstacles to its use in
Scotland. First, there are several technical obstacles:
\begin{itemize}
\item
  Several of the links used by Tegola and related projects are
  significantly longer 
  than 13km.
\item
  Transmission in higher frequencies is adversely affected by high 
  humidity and high temperatures. Typically, only one of 
  these is a significant problem in Scotland..
\item
  The Ubiquiti equipment uses substantially more power than their 5GHz 
  offerings -- about 40W. This would make it unsuitable for solar and 
  wind-powered relays.
\end{itemize}

Second, there is a regulatory obstacle.
Ubiquity claim that, \emph{The airFiber®24 operates in worldwide,
  license-free, 24 GHz frequencies. Users can deploy airFiber24 almost
  anywhere they choose (subject to local country regulations).}
  In the UK, the 24.05--24.25GHz band used by this equipment is
  allocated by Ofcom to radiolocation (1) and amateur radio
  (2).\footnote{\url{http://www.ofcom.org.uk/static/archive/ra/publication/ra_info/ra365.htm}} Their
  regulation includes the following conditions:
\begin{enumerate}
\item Non-government low power devices in the radiolocation services
  are limited to:
\begin{enumerate}
\item Portable and fixed applications between 24.15-24.25 GHz; and
\item Mobile applications between 24.25-24.35 GHz on a NIB to the radionavigation service
\end{enumerate}
Power flux-density at 10 metres from the system antenna in the
direction of maximum radiation is not to exceed 1.5 mW/m2 with out
approval
\item 24.05-24.25 GHz is used by the Amateur service. The part of the allocation between 24.05 and 24.150 GHz may only be used with written consent of the Secretary of State
Home Office/Office of The Scottish Executive for the Emergency Services between 24.05-24.15 GHz
ISM apparatus may use the band 24.0-24.25 GHz.
\end{enumerate}
The 24.05--24.35GHz band band is partitioned into three
  sub-bands, two of which fall within the 24.05--24.25GHz range of the
  Ubiquity equipment. Of these two, one
  is reserved for government and amateur radio use and is not allowed
  for general use and the last is permitted at extremely low power
  densities for radiolocation.  

We obtained a ``non-operational'' licence from Ofcom in order to test the equipment at its rated
  power of  $13\,\text{W}/\text{m}^2$ and in both sub-bands. Operational
  use of this equipment for community broadband would require
  regulatory changes.



\subsection{Link budgets}

This experiment is meant to compare the real world performance against
the vendor's claims, but we first considered how the claims stack up
against the theoretical performance. We planned to run this equipment
on a link of around
$6\,\text{km}$. Our calculations suggested that
using $40\,\text{MHz}$ wide channels we should be able to solidly
sustain such links at $256\,\text{Mbps}$ or better for shorter links.

Our initial plan was to test the equipment on existing Tegola
relays one is a 6.5km link; the second 15.5km. Although the latter is
over the advertised range, even a substantial fraction of the
advertised throughput would be useful.

The following is a roughly chronological account of the project.
The initial installation was done during a period of very high winds
in early January 2014.

\subsection{Initial configuration and testing}
\label{december-2013-initial-configuration-and-testing}

\begin{wrapfigure}{r}{0.3\textwidth}
\includegraphics[width=0.3\textwidth]{radio-in-corridor}
\end{wrapfigure}
We ordered and received one pair of radios. Before deploying them
we thought it would be a good idea to check that they were working
and test them in ideal situation -- our office corridor. One thing
we immediately noticed was how critical alignment is. Even over
a distance of 35m, the performance fell of dramatically if the
antennae were slightly out of alignment. It's a very good idea to
configure equipment before deploying it, but to do this we had to turn
off sychronisation which relies on GPS and doesn't work indoors.

\subsection{Strengthening the masts}
\label{december-2013-january-2014-strengthening-the-relays}

Our basic relay construction (see the
\href{http://www.tegola.org.uk/howto/relay-construction.html}{relevant
  howto on the Tegola web site}) uses aluminium pegs to anchor the
diagonal braces to the ground. Both sites were on terrain that
consisted of bedrock covered by peat of varying depth. Although we
have never had a problem with the pegs shifting, peat is a bit
jelly-like, and the structures can wobble through a cm or two. The
alignment of 24GHz is much more critical than for the lower bandwidths
of 2.4 and 5.8GHz, so we replaced the pegs with epoxy bolts into the
bedrock.
\begin{figure}[h]
\includegraphics[width=0.4\textwidth]{corran-peg-s}
\includegraphics[width=0.4\textwidth]{corran-epoxy-s}
\caption{Pegged anchors (left) and epoxied anchors (right).}
\end{figure}

We also strengthened both relays. For example, at one of our relays we
added an extra horizontal bar.

\begin{figure}[h]
\includegraphics[width=0.4\textwidth]{corran-before-from-behind-s}
\includegraphics[width=0.4\textwidth]{corran-after-from-behind-s}
\caption{Corran mast before (left) and after (right)}
\end{figure}

\subsection{Installation and alignment}
\label{january-2014.-installation-and-alignment}

The radios are reasonably light (under 10kg) but awkward to carry up
hills. We used an old backpack frame that 30 years ago had been used
for carrying batteries up to community TV relays. The radios come with
a mounting frame that is first attached to the structure. The radio is
then ``slotted'' into the mounting frame. This arrangement makes it
quite easy to install the whole assembly when working from a ladder.

The antenna can be aligned through elevation and azimuth
adjustment screws. Unfortunately there is a great deal of backlash in
these screws, and they are almost useless if you are working in high
winds. If the clamping bolts are loose, the antenna is blown around
through the considerable travel allowed by the adjustment screws. The
signal strength read-out is at the bottom of the antenna, and if
the alignment is being done from a ladder, you almost certainly
need someone below (with a hard hat) to squint up and call out
the figures.

The installation instructions recommend an alternating process in which
one end of the link is adjusted then the other, and so
on. Unfortunately we were unable to complete this process before
the weather closed in and our workforce departed. However, the
alignment is good enough that we can start taking some measurements.
The initial indications are that the link will work reasonably well
over a distance of 6.5km.

\subsection{Performance}
\label{performance}

Overall performance of the 24GHz was quite good with throughput of
~500Mbps easily achievable over the distance from Corran to Coille
Mhialairigh, a distance of about 6.5km. This is more than was easily
available at the time in the the 5GHz band. It is also less than
advertised, which is perhaps unsurprising since the headline figures
reflect what is possible under ideal conditions.

\begin{figure}
\includegraphics[width=0.6\textwidth]{corran-rcvcap}
\caption{Corran receive capacity}
\label{fig:corran-rcvcap}
\end{figure}
Figure~\ref{fig:corran-rcvcap} shows the link speed as measured at
Corran for a typical day. The fluctuations are not obviously
correlated with any specific phenomenon and a link budget calculation
compared with the specification sheet for which
modulation schemes are available at which signal and noise levels
suggests that the capacity negotiation is overly aggressive. That is
to say the radios are trying to maintain a link speed somewhat higher
than their spec sheet says they should. Sometimes they actually
achieve this, but unstably.

What we can conclude from this is that for the link to be stable, the
modulation scheme must be explicitly set to one that can easily be
supported by the link with a sufficient margin and not allow the
automatic negotiation to proceed. Incorporating such a margin (we
recommend at least 6dB over the \emph{observed} signal levels)
provides a slower but much more stable link.

\subsection{5GHz Tidal Fading}
Tidal fading has been a serious problem in over-water links that 
use to the 2.4 and 5.8 GHz ranges.  They are less of a problem with
higher frequencies (e.g. 24GHz) in which the Fresnel ellipsoids and
beamwidth are narrower~\footnote{our 24GHz link is not over
  water}. The 5GHz airFiber link operates over a distance of
approximately  20.8km with the two ends at heights of 300m and 20m
above MSL.  We have seen some tidal fading on normal 5GHz transmission
over this link and disruptive fading over the same stretch of water
when both ends are at heights of approximately 20m.

\begin{figure}
\includegraphics[width=0.7\textwidth]{tidedata/rxpower0.jpg}
\includegraphics[width=0.7\textwidth]{tidedata/rxpower1.jpg}
\caption{Scatter plot of signal strength against sea level for
  horizontal (top) and vertical  polarizations}
\label{fig:rxpower}
\end{figure}

In January 2015 we had a failure of our the operating link and, since
the weather prohibited climbing up to the 300m high relay, we switched
to the experimental 5GHz airFiber equipment.
Figure~\ref{fig:rxpower} shows the received signal strengths on both
polarizations.  The classic quasi-cycloidal pattern is discernable in
both plots, however it is in a range over with dips to -78 dBm, which
might cause problems.  However these extreme points are probably due
to other causes. Examination of the RX capacity (see
Figure~\ref{fig:capacity} indicates that other effects may dominate tidal
fading and that other things being equal, tidal fading is likely to be
less of a problem on these links than it is on conventional 5GHz
equipment. Readings were taken at 10 minute intervals from mid-January
to mid-February
\begin{figure}
\includegraphics[width=0.7\textwidth]{tidedata/rxcapacity.jpg}
\caption{Scatter plot of RX channel capacity against sea level for
  horizontal (left) and vertical  polarizations}
\label{fig:capacity}
\end{figure}

\begin{wrapfigure}{r}{0.3\textwidth}
  \includegraphics[width=0.3\textwidth]{tommy-diversity}
\end{wrapfigure}
Tidal fading, or tidal multi-path interference is a known problem in
radiofrequency engineering. It happens when the direct signal from
a transmitting antenna and an indirect or reflected signal both arrive at
the receiving antenna with opposite phase and cancel each other
out. The degree to which this happens depends on how long the
reflected path is, the direct path being fixed. Since in this case the
reflections are caused by the water, as the tide rises and falls the
reflected path gets longer and shorter. At certain points in the
water's travel, the signals cancel out nearly completely -- from the
plots the loss from this effect is 12dB or a signal around 16 weaker
than it would normally be.

The solution is to have multiple antennae spaced so that when
the signal reaching one is being destructively interfered with, the
other has a good signal and vice-versa. The picture to the right shows
an inexpensive arrangement that implements this. It uses the standard
5GHz rocket radios from Ubiquiti and a non-standard arrangement with
two aftermarket antennae. This has been found to work very
well. Unfortunately such a solution is not possible with the AirFiber
units because of their construction -- the radios and the antennae are
integrated into a single housing, making it impossible to use
different antennae. Nevertheless the AirFibre would be useful for
links that do not traverse tidal bodies of water.

\clearpage
\section{Free-space Optics}
\label{sec:fso}



At the request of the Scottish Government we also evaluated an
unrelated technology for use in urban areas – free-space optics. This
means signalling by laser through the air instead of through
fibre-optic cabling. The main question here was, since these devices
operate in the visible spectrum, how well do they operate in reduced
visibility conditions? Do they operate well enough for use in the Scottish
climate over short distances. Again it is a question of cost. The
equipment is expensive (some \pounds 8-15k per link) and requires careful
mounting and alignment, but compared to the civil engineering cost of
running fibre in urban areas, or leasing it, it may be worth it.

The free-space optics equipment, as expected, performs well over
distances of a few hundred metres, but will not be useful for longer
links in a typical Scottish environment. The equipment is expensive,
and beyond the reach of a small community project, though we might
have been able to reduce the cost by using an alternative vendor.

A link was established between TechCube and Appleton Tower, a distance
of 500m, and evaluated using routine HUBS traffic.

Installing the link was significantly more arduous than installing the
accompanying wireless link. It required on-site optical splicing to
bring the optical link to the server cabinets at each end-point, and a
dedicated high-current power supply.

It was also suggested that custom arrangements, more rigid than those
required for a wireless connection, would be required for a laser
link. In the event, we used some of our standard fixing arrangements
(bolting to a framework of alloy scaffold poles anchored to the
building at 3m intervals) for the Summerhall end of the link, and this
was unproblematic.

The purpose of the trials with the free-space optical (FSO) links was to
ascertain the extent to which fog, a common feature of the Scottish
environment especially in coastal areas, would affect their
use. Aerosolised water droplets attenuate visible and infrared
light. We can quantify this effect theoretically, but we were
primarily interested in finding out if their operational window was,
on average, large enough to justify the considerable expense.
\begin{figure}
  \includegraphics[width=\textwidth]{at-laser-s.jpg}
\end{figure}

We found that performance of the link correlates well with standard
MIDAS visibility records at a nearby weather station, and recommend
that historical visibility data should be analysed to assess the extent
to which fog and rain will affect the performance of such a link in
any proposed location.

\subsection{Link Design}
\label{sec:link-design}

CableFree, the laser vendor produced a link design as shown in
Figure~\ref{fig:cablefree-link}. In principle it is plausible,
consisting of a laser link, and a back-up radio link between Appleton
Tower and the Tech Cube, a distance of about 500m. At either end two
switches would use the spanning-tree protocol (STP) configured with the
radio link having a higher cost than the laser link. The principle
being that traffic would flow over the lasers unless they were
unavailable in which case the backup link would be used.
\def\cfdesign{%
    \node[draw,rotate=90] (atnetgear) at (0,0) {Netgear Switch};
    \node[draw] (atlaser) at (2,1) {Laser};
    \node[draw] (atradio) at (2,-1) {Radio};
    \node[draw] (tclaser) at (6,1) {Laser};
    \node[draw] (tcradio) at (6,-1) {Radio};
    \node[draw,rotate=-90] (tcnetgear) at (8,0) {Netgear Switch};
    \draw[thick,blue] (-1,0) edge (atnetgear.north);
    \draw[thick,blue] (9,0) edge (tcnetgear.north);
    \draw[thick,orange] (atnetgear.south) edge (atlaser.west);
    \draw[thick,orange] (tcnetgear.south) edge (tclaser.east);
    \draw[thick,blue] (atnetgear.south) edge (atradio.west);
    \draw[thick,blue] (tcnetgear.south) edge (tcradio.east);
    \draw[thick,red,dotted] (atlaser.east) edge (tclaser.west);
    \draw[thick,black,dotted] (atradio.east) edge (tcradio.west);
    \node at(4,0) {$\approx 500\text{m}$};
    \node at(4,1.5) {$\text{cost} = 1000$};
    \node at(4,-1.5) {$\text{cost} = 100000$};
}
\begin{figure}[h]
  \centering
  \begin{tikzpicture}
    \node at (1,2) {\textbf{Appleton Tower}};
    \node at (7,2) {\textbf{Tech Cube}};
    \cfdesign
  \end{tikzpicture}
  \caption{CableFree Link Design}
  \label{fig:cablefree-link}
\end{figure}

We already had a significant amount of radio equipment at
both sites -- including a radio link between them -- so an extra radio
link could be considered superfluous. 
\begin{wrapfigure}{r}{0.3\textwidth}
  \includegraphics[angle=-90,width=0.3\textwidth]{also-faulty-s}
\end{wrapfigure}
The mounting brackets supplied with these back-up radios were somewhat flimsy. It can get very windy at the top
of tall buildings in Edinburgh. Wind speeds can reach several times
more than at ground level. The installation is shown at right with the
mounting bracket having worked loose in high winds leaving the radio
pointing downwards. 

We also had high end HP ProCurve J9050A switches at either end (loaned
by the School of Informatics), far more capable network elements than
the consumer-grade Netgear switches (see
Figure~\ref{fig:carrier-grade}) specified by CableFree. The
deficiencies of the Netgear switch are that it is not manageable via
telnet or ssh, which makes it very difficult to recover from outages
or erroneous configurations in a network of any size or complexity.

 We ultimately ended up operating the link by
connecting the vendor-supplied netgear switches to our HP switching
core and removing the vendor-supplied radio equipment entirely.
\begin{figure}
  \begin{center}
    \includegraphics[width=0.5\textwidth]{carrier-grade-s}
  \end{center}
  \caption{A ``carrier-grade'' ethernet switch.}
  \label{fig:carrier-grade}
\end{figure}

The more serious deficiency with this design is the use of the
spanning-tree protocol. In brief, the mechanism is unstable in marginal
conditions. If the weather is clear or if the weather is very foggy,
the link operates as designed. However it appears that,in the region between clear and
very foggy, the path is liable to rapidly change between the (possibly
unuseable) optical link and the radio link. 



\subsection{Installation}
\label{sec:fittings}

\begin{wrapfigure}{r}{0.3\textwidth}
  \includegraphics[angle=-90,width=0.3\textwidth]{tada-s}
\end{wrapfigure}
Two different methodologies were used for the physical
installation. At Appleton Tower we were not permitted to make any
permanent modifications to the building. The laser head was therefore
mounted on a pedestal, shown at right. This was fixed to a paving tile
and further weighed down with more scrap paving tiles to prevent it
from blowing away.

Installation at the Summerhall site was straightforward since we were
able to make building modifications. We used aluminium scaffolding to
create a frame fixed to the side of the plant room on the roof and
fixed the laser head to that.

Alignment of the lasers was straightforward -- in fact easier than
anticipated and though it was done with specialised equipment (which
measures signal strength and produces audible faster or slower beeps
accordingly as the optical link is well or poorly aligned) it appears
to be quite possible to align the link without any more equipment than
a simple two-way radio. This is done by observing the beam edges on
one end (after a rough alignment) and instructing a colleague over the
radio to slightly move the far laser until the edges are centered with
respect to the near one, and repeating the procedure for the other
side. This turned out to be necessary on one occasion, as severe winds caused the
Appleton Tower pedestal-mounted end to de-align.

We also had to make some modifications to the lasers due to poor
design. A protective piece of glass that covers the lens had to be
fixed with screws because it came off easily. Luckily we found this,
to our surprise, in the workshop before installing the equipment.

\begin{figure}[h]
  \begin{center}
    \includegraphics[width=0.5\textwidth]{faulty}
  \end{center}
\end{figure}



\subsection{Instrumentation}

The laser and wireless links were instrumented using standard port
statistics on the HP switches, collected in the usual way at five
minute intervals. It would have been interesting to be able to collect
more detailed information about the link -- signal strength, bit-error
rate and so forth -- but this was not possible as the vendor required
executing an unacceptable non-disclosure agreement in order to be able to access this basic information and this would have prevented
dissemination of these results.

The port statistics are stored in the industry-standard Round-Robin
Archive format.


\subsection{Correlation of link operation with weather}
\label{sec:weather}

We have been monitoring availability of the optical link and we have
access to the UK hourly weather data for the purposes of this project
Unfortunately we had to undertake some emergency works in
November to make way for Vodaphone at Summerhall, and it turns out
that this interfered with some of our data-collection, and we were
unable to make planned fine-grain observations that might have given
more detail on the inferred instabilities of the spanning-tree algorithm.


\begin{figure}
  \centering
  \includegraphics[width=\textwidth]{results/fig1-Laser-12-Apr-12-Nov}
  \caption{Laser}
  \label{fig:throughput_laser}
\end{figure}
Figure~\ref{fig:throughput_laser} shows the daily throughput in either
direction of the laser link over an eight-month period (blue and green
curves). The daily volume of traffic grows fairly steadily over the
eight-month period. These curves are overlaid on the hourly record of
visibility at Gogarburn (the closest MIDAS station for which hourly
visibility records are available), for reasons that will soon be made
clear.

The weather data shows no correlation between wind-speed and quality
of the FSO connection (we had conjectured that there might be some
effect if the wind routinely affected our mountings and beam alignment). Both
Summerhall and the Appleton tower are relatively low-rise
and solid buildings. On high-rise buildings wind-induced movements
could have a significant effect (for example, in high wind the Shard
is intended to move by up to 20 inches, the Burj Khalifa moves by 2 metres). The beam
was observed to have an effective
spread of around 2m (diameter) over the 500m distance (around 13’ of
arc) and it was relatively straightforward to establish a functional
alignment, which was maintained without further intervention over the
course of this experiment. 

However, we did observe, as anticipated,
that the FSO link became effectively inoperative in conditions of
severely reduced visibility.

\begin{figure}
  \centering
  \includegraphics[width=\textwidth]{results/fig2-Wireless-12-Apr-12-Nov}
  \caption{Wireless}
  \label{fig:throughput_wireless}
\end{figure}
The radio link is required to maintain service when the FSO is
inoperative. It comes into play automatically when FSO performance is
degraded to such an extent that the radio link is better. Figure 2
shows the daily throughput of the radio link, again overlaid on
visibility data.

\begin{figure}
  \centering
  \includegraphics[width=\textwidth]{results/fig3-Wireless-28-March-1-June}
  \includegraphics[width=\textwidth]{results/fig4-Wireless-14-July-14-Aug}
  \includegraphics[width=\textwidth]{results/fig5-Wireless-September}
  \caption{Detail}
  \label{fig:throughput_detail}
\end{figure}
Figure~\ref{fig:throughput_detail} show portions of this period in
more detail. They show that use of the radio link is directly
correlated with episodes of low visibility (below 500m). Going back to
figure 1 we see that many of these episodes are short-lived, and have
little effect on the overall daily volume of traffic, but that in some
cases the fog is more persistent.

To quantify the likely impact of these climatic effects in any
proposed location, it would be useful to study both the frequency and
the duration of local episodes of low visibility. Longer episodes will
have greater impact in many applications.

We also note that longer links will clearly be more susceptible to
such effects. For example, over the course of 2014, only 10\% of hourly
observations at Gogarburn recorded a visibility of 500m or less, while
the figures for limits of 1km, 1.5km, and 2km are 10\%, 18\% and 44\%
respectively.

Unfortunately, as mentioned above, we were unable to monitor the
performance of the FSO link directly because of the proprietary
protocol that prevented adequate instrumentation. Nevertheless, we
were able to quantify the difference in performance between the laser
link and a back-up wireless link running over 5GHz unlicensed
spectrum. The FSO link has several times the bandwidth of the radio
link, but it is more susceptible to climatic degradation.

%\printbibliography
\end{document}

%%% Local Variables:
%%% mode: latex
%%% TeX-master: t
%%% reftex-default-bibliography: ("literature.bib")
%%% zotero-collection: #("28" 0 2 (name "Papers/Scotgov Report"))
%%%